\documentclass[10pt,a4paper]{amsart}
\usepackage[margin=19mm]{geometry}
\usepackage{amsmath,amssymb,mathtools}
\usepackage[scr=rsfs,cal=euler]{mathalfa}
\usepackage{xfrac}
\usepackage[unicode,hidelinks,bookmarks=false]{hyperref}
\newcommand{\ie}{i.e.\ }
\newcommand{\cf}{cf.\ }
\newcommand{\resp}{respectively\ }
\newcommand{\Subsection}{\subsection}
\providecommand{\nobreakdash}{\nobreak\hbox{-}}
\setlength{\emergencystretch}{2em}
\multlinegap0pt
\usepackage{amssymb}
\usepackage{bm}
\usepackage{mathtools}
\usepackage{xcolor}
\newtheorem{thm}{Theorem}
  \newcommand{\KIDi}{\operatorname{KID}^{(1)}}
  \newcommand{\KIDii}{\operatorname{KID}^{(2)}}
  \newcommand{\uKIDi}{\operatorname{uKID}^{(1)}}
  \newcommand{\uKIDii}{\operatorname{uKID}^{(2)}}
  \newcommand{\uKIDo}{\operatorname{uKID}^{(0)}}
\title{Unit Killing Initial Data}
\author{Igor Khavkine$^\flat$ {} and Salvador Mengual Sendra$^{\sharp}$}
\date{Source: arXiv:2607.21815v1 (CC BY 4.0)}
\begin{document}
\maketitle

\noindent\emph{Source and licence.} Unit Killing Initial Data. Version arXiv:2607.21815v1 (CC BY 4.0), distributed by arXiv under the Creative Commons Attribution 4.0 International licence, http://creativecommons.org/licenses/by/4.0/, as recorded in the arXiv licence metadata for this exact version and shown on the abstract page https://arxiv.org/abs/2607.21815v1. Full licence notice: \texttt{licenses/arxiv-2607.21815-CC-BY-4.0.txt}.\par\smallskip

\noindent\textit{Editorial scope.} Counted unit: the article's thm thm:canuKID (source0.tex lines 552--633). The article's complete original proof of it is delivered with it (source0.tex lines 635--679, one \texttt{proof} environment of 45 lines). No mathematics is written, repaired or re-derived: the statement and the proof are the article's own lines, in the article's own notation, and no retained line is substituted or re-flowed.  Retained uncounted as the source-local context the proof needs: lines 408--418; lines 520--524.  Nothing was omitted from any retained span, so the package carries no omission record.  No retained line contains a citation, so the wrapper carries no bibliography and no bibliography entry.  No retained line contains a figure environment, an image-inclusion command or drawing code, so no image is delivered and the package declares no asset dependency.  Editorial formatting only, in addition: an AMS \texttt{amsart} preamble that restates the article's own declarations and macros used by the retained lines, namely the environments \texttt{thm} on separate counters, together with the standard packages those lines need.  The retained environments print in this document as Theorem 1 in order of appearance, whereas the article prints the same items with the numbering of its own full document.  The complete native source of the article as distributed in the arXiv e-print for this exact version is bundled unchanged as \texttt{sources/205894/source0.tex}.
\begin{equation} \label{eq:KID}
\begin{aligned}
	\KIDi_{AB}[v] &:= D_A v_B + D_B v_A + 2\pi_{AB} v_0 = 0 ,
	\\
	\KIDii_{AB}[v] &:= D_{A} D_{B} v_{0}
	+ (2(\pi\cdot\pi)_{A B} + \nabla_0 \pi_{AB}) v_0 \\
	& \quad {}
	+ 2 \pi_{(B}{}^{C} D_{A)} v_C
	+ (D^{C}{\pi_{A B}}) v_{C} = 0 ,
\end{aligned}
\end{equation}

\begin{equation} \label{eq:uKID}
	\uKIDo_{AB}[u] = 0 , \quad
	\uKIDi_{AB}[u] = 0 , \quad
	\uKIDii_{AB}[u] = 0
\end{equation}

\begin{thm}[prolongued KID and uKID equations] \label{thm:canuKID}
(a) The following system is equivalent to \eqref{eq:KID}:
\begin{equation} \label{eq:canKID}
\begin{aligned}
D_A v_B &= -\pi_{AB} v_0 + \omega_{AB} , \\
D_A v_0 &= p_A , \\
D_A \omega_{BC} &= 2 (\pi_{A[B} p_{C]} - v_0 D_{[B} \pi_{C]A}) - r_{BCAD}
u^D , \\
D_A p_B &= -2\pi_{(A}{}^C \omega_{B)C} - u^C D_C \pi_{AB} - v_0 \nabla_0 \pi_{AB} .
\end{aligned}
\end{equation}
The original equations on $v_A$ are recovered by using the following
substitutions, obtained from the first two equations above:
\begin{equation} \label{eq:canKID-defs}
	\omega_{AB} = D_{[A} v_{B]} , \quad
	p_A = D_A v_0 .
\end{equation}

(b) The following system is equivalent to \eqref{eq:uKID}:
\begin{equation} \label{eq:canuKID}
\begin{aligned}
	D_A u_B &=
		\omega_{AB} - \pi_{AB}\,u_0 + 2 q_{(A} u_{B)},
\\
	D_A q_B &=
		\frac12 \omega_A{}^{C}\omega_{BC} + \frac52 q_A q_B + \omega_{C(A} u_{B)} q^{C} + \frac12 q_C q^{C} u_A u_B 
\\
& \quad
		+ \frac12 (\pi_{AC} \pi_{BD} - \pi_{AB} \pi_{CD} - r_{ACBD} ) u^{C} u^{D} - \frac{\pi_{AB} q_C u^{C}}{u_0}
\\
& \quad
		- \frac{1}{2 u_0^2} (q_A q_B + q^{C} q^{D} u_A u_B u_C u_D + 2\omega_{D(A} q^{C} u_{B)} u_C u^{D}
\\
& \quad
		- 2q^{C} u_C q_{(A} u_{B)} - 2\omega_{C(A} q_{B)} u^{C} + \omega_{AC} \omega_{BD} u^{C} u^{D} )		
\\
& \quad
		- u_0u^{C} ( D_{(A} \pi_{B)C} + D_C \pi_{AB} )
\\
& \quad
		+ \frac12 u_0^2 (\nabla_0 \pi_{AB} + \pi_A{}^{C} \pi_{BC}),
\\
	D_A \omega_{BC} & = 
		2 \omega_{BC} q_A + 2 \omega_{A[C} q_{B]} + \omega_A{}^D \omega_{D[C} u_{B]} - 2 q_A q_{[C} u_{B]}
\\
& \quad
		 + \pi_{A[C} \pi_{B]D} u^D - r_{BCAD} u^D + r_{ADE[B} u_{C]} u^D u^E
\\
& \quad		
		+ \pi_{DE} \pi_{A[C} u_{B]} u^D u^E + q^D \omega_{D[C} u_{B]} u_A
\\
& \quad
		+ \pi_{AD} \pi_{E[B} u_{C]} u^D u^E - 2 \pi_{A[C} q_{B]}
\\
& \quad
		- \frac{1}{u_0} (2 \pi_{A[C} \omega_{B]D} u^D - 2 \pi_{A[C} q_{B]} + q_A q_{[C} u_{B]})
\\
& \quad
		+ \frac{1}{u_0^2}(\omega_{D[C} u_{B]} q_A u^D - \omega_{AD} q_{[C} u_{B]} u^D - q^D q_{[C} u_{B]} u_A u_D 
\\
& \quad	
		+ q_A q_{[C} u_{B]} + q^D\omega_{E[B} u_{C]} u_A u_D u^E - \omega_{AD} \omega_{E[C} u_{B]} u^D u^E)  
\\
& \quad
		- u_0 [ (D_A \pi_{D[B}) u_{C]} u^D + 2D_{[B} \pi_{C]A} + u^D u_{[C} (D_{B]} \pi_{AD}) 
\\
& \quad
		+ 2 (D_D \pi_{A[C}) u_{B]} u^D + 2 \pi_{A[C} q_{B]}]
\\
& \quad
		- u_0^2 (\nabla_0 \pi_{A[C} + \pi_A{}^D \pi_{D[C}) u_{B]}.
\end{aligned}
\end{equation}
The original equations on $u_A$ are recovered using the following
substitutions, obtained from the first equation above:
\begin{equation} \label{eq:canuKID-defs}
\begin{aligned}
	\omega_{AB} &:= D_{[A}u_{B]}, \\
	q_A &:= \left( \frac{\pi_{AB} u^B}{u^D u_D} - \frac{\pi_{BC} u^B u^C}{2(u^D u_D)^2} u_A \right) u_0 + \frac{u^B(D_{(A} u_{B)})}{u^D u_D} - \frac{u^B u^C (D_C u_B)}{2(u^D u_D)^2} u_A .
\end{aligned}
\end{equation}
\end{thm}

\begin{proof}
(a) One direction is straightforward. Given a solution
of~\eqref{eq:canKID}, its $v_0$ and $v_B$ components solve the KID
system~\eqref{eq:KID}, which can be checked by mechanically
substituting~\eqref{eq:canKID} into~\eqref{eq:KID} and eliminating
$\omega_{AB}$ and $p_A$ using the definitions~\eqref{eq:canKID-defs}. A
quick inspection of definitions~\eqref{eq:canKID-defs} shows that they
themselves follow from~\eqref{eq:canKID}, so they are not separate
equations. In the other direction, given a solution $v_0$ and $v_B$ of
the KID system, it solves the prolonged system~\eqref{eq:canKID} when
the extra $\omega_{AB}$ and $p_A$ components are defined
by~\eqref{eq:canKID-defs}. Up to using these definitions, the $D_A v_B$
equation coincides with $\frac{1}{2} \KIDi_{AB}[v]$, the $D_A v_0$
equation is trivial, the $D_A \omega_{BC}$ equation coincides with
$D_{[B} \KIDi_{C]A}[v]$, while the $D_A p_B$ equation coincides with
$\KIDii_{AB}[v] - \pi_{(A}{}^C \KIDi_{B)C}[v]$. These maps clearly give
a bijection on solutions.

(b) The proof is essentially analogous to part (a), though
computationally more involved. The straightforward direction and the
bijectivity follow the same logic. In the other direction, we get a
solution of the prolonged system~\eqref{eq:canuKID} from a solution of
the uKID system~\eqref{eq:uKID} when the extra $\omega_{AB}$ and $q_A$
components are defined by~\eqref{eq:canuKID-defs}. Up to using these
definitions
\begin{align*}
	(D_A u_B + \cdots)
		&= \uKIDi_{AB}[u] , \\
	(D_{[A} q_{B]} + \cdots)
		&= \frac{1}{4} \uKIDo_{AB}[u] + O(\uKIDi[u]) , \\
	\MoveEqLeft
	(D_A \omega_{BC} - 2 u_{A} D_{[B} q_{C]} + 2 u_{[B} D_{C]} q_{A} + \cdots)
		\\
		&= 4 D_{[C} \omega_{B]A}  - 3 D_{[C} \omega_{BA]} + O(\uKIDi[u]) , \\
	\MoveEqLeft
	(D_{(A} q_{B)} + \frac{1}{2} (u_{A} u^C D_{[B} q_{C]} + u_{B} u^C D_{[A} q_{C]}) + \frac{1}{2} u^C D_{(A} \omega_{B)C} + \cdots)
		\\
		&= \frac{u_0}{2} \uKIDii_{AB}[u] + O(\uKIDi[u]) ,
\end{align*}
where $\cdots$ stand for terms of order zero and the omitted terms on
the right-hand side that are not explicitly shown are long and otherwise
unenlightening. From the structure of the left-hand sides above, it is
clear how to solve for $D_A u_B$, $D_A q_B$ and $D_A \omega_{BC}$ by
algebraic operations.
\end{proof}

\end{document}
